\scanpage{90}
\mspage{44}
10) \uline{Structure} MDT \uline{discrète} : \uline{action} des \uline{morphismes de normalisation entre \inserted{MD-}graphes}.

Nous voulons exprimer, en plus des morphismes \deleted{$\varphi_G$ et} $S\Pi f$ ($f$ morphisme de générisation) et $\Pi f$ ($f$ iso de MD-graphes) entre autres \unsure{des|ces} morphismes du type
\[
\diagram{\begin{tikzcd}[ampersand replacement=\&]
M_G \arrow[r] \& M_{g_\alpha,\widehat I_\alpha}
\end{tikzcd}}
\]
et
\[
\diagram{\begin{tikzcd}[ampersand replacement=\&]
M^{\unsure{?|!}}_{g,\unsure{I|\overline I}} \arrow[r] \& M_{g,I'}
\end{tikzcd}}
\qquad\text{pour }I'\subset I.
\]
On va introduire pour ceci une autre notion de morphismes de MD-graphes, qu'on appellera morphismes de normalisation (l'intersection des deux notions de morphisme est celle d'iso de MD-graphes).

Géométriquement, cela correspondra à la situation d'une courbe stable \deleted{$(\widehat X,I)$}\inserted{$X=\widehat X\setminus I$} sur un corps alg. clos $\unsure{k|h}$, en distinguant une partie $I'\subset I$ de l'ens. des pts à l'infini, et une partie $A'\subset A$ de l'ens. des pts doubles de $X$,
\margin{89}
